\section{Method}
\label{sec:method}

\subsection{TIGER Fingerprint Space}

We embed all 150 corpus games in the five-dimensional TIGER space using the
within-game z-score normalization and PCA projection described in TG-A2 and
TG-B2~\cite{dellaLibera2026pca, dellaLibera2026fingerprint}. Each game is represented
as a five-dimensional TIGER coordinate $(T, I, G, E, R)$ where each dimension is
on a 0--10 scale.

\subsection{Sparse Cluster Identification}

K-means clustering with $k = 6$ (selected by elbow method on within-cluster inertia
over $k \in \{3, \ldots, 10\}$) is applied to the TIGER coordinate space.
Cluster density is measured as the number of corpus games per unit volume of
TIGER space (estimated via Gaussian kernel density estimation with bandwidth 1.5 in
standardised TIGER coordinates). Cluster centroids with density in the bottom quartile
are classified as sparse, and sparse coordinates at maximum distance from all corpus games
are selected as design gap candidates.

\subsection{Expert Validation Study}

Sparse cluster descriptions are translated into natural-language design briefs:
one paragraph describing the design character implied by the cluster's TIGER
profile, without reference to any existing game. Five experienced game designers
(each with $\geq 2$ published titles) participate in a structured interview:
they read each brief, rate its commercial viability (1--5 scale), and assess
whether any existing game already occupies that space. Briefs rated $\geq 3.5$
by $\geq 3$ of 5 designers and assessed as genuinely unexplored by $\geq 3$
designers are classified as validated opportunities.
